\section{A Hilbert kernel with chart extensions}
\label{ker:sec}

Retain the directions, maps, and pseudoinverses from
\Cref{lem:directions}. We construct common Hilbert features on the parameter
space and extend them to the rounded chart labels. Their squared distances
will approximate Euclidean distance. We will control violations of the
triangle inequalities both globally and in proportion to within-chart
coordinate distance. These are the inputs for the metric construction in
\Cref{sec:metric}.

The parameters in this section are
\begin{equation}
\label{ker:parameters}
 a=m^{-40},\qquad b=m^{40},\qquad
 \varepsilon=m^{-200},\qquad r=m^{-8},\qquad \tau=m^{-2000}.
\end{equation}
The numerical exponents only enforce separation of scales. The cutoffs
$a,b$ approximate Euclidean distance; $\varepsilon$ controls the added
Fourier components; $r$ is the tolerance for the angle repair; and $\tau$
is small enough to absorb every polynomial norm bound below. All choices
are fixed powers of $m$, which is essential for the vertex count.

\subsection{The common kernel}

The Gaussian block below produces the macroscopic Euclidean distance
by integrating positive-definite Gaussian kernels over scales, as in the
classical negative-type framework of Schoenberg~\cite{Sch38}.
The additional Fourier blocks have negligible distance contribution but
will supply the coordinate profiles for the chart extensions.

Write $E(t)=(\cos t,\sin t)$, and let $\gamma_m$ be standard Gaussian measure
on $\R^m$.  Define a map $p:\R^m\to\Hh$ into the real Hilbert space
\[
 \Hh=
 L_2([a,b]\times\R^m,d\sigma\,d\gamma_m;\R^2)
 \mathbin{\oplus}
 \bigoplus_{s=1}^{S}\bigoplus_{i=1}^{N}
 L_2([a,b],d\sigma/\sigma;\R^2)
\]
by assigning to $p(\theta)$ the first component
$m^{1/4}E(\langle g,\theta\rangle/\sigma)$ and the component
$\sqrt\varepsilon E(\langle g_{s,i},\theta\rangle/\sigma)$ in summand
$(s,i)$.  The Gaussian characteristic function gives
\begin{equation}
\label{eq:kernel}
\begin{split}
 \langle p(\theta),p(\eta)\rangle
 &=\sqrt m\int_a^b
     e^{-|\eta-\theta|^2/(2\sigma^2)}\,d\sigma\\
 &\quad+\varepsilon\sum_{s=1}^{S}\sum_{i=1}^{N}
   \int_a^b
   \cos\!\left(\frac{\langle g_{s,i},\eta-\theta\rangle}{\sigma}\right)
   \frac{d\sigma}{\sigma}.
\end{split}
\end{equation}

\begin{lemma}[Approximation of Euclidean distance]
\label{lem:kernel}
There is an absolute constant
\[
 c_*=2\int_0^\infty\bigl(1-e^{-1/(2t^2)}\bigr)\,dt>0
\]
such that, whenever $|\theta|,|\eta|\le3\sqrt m$,
\begin{equation}
\label{ker:distance}
 \left|\,|p(\theta)-p(\eta)|^2
       -c_*\sqrt m\,|\theta-\eta|\,\right|\le m^{-37}.
\end{equation}
\end{lemma}

\begin{proof}
The integral defining $c_*$ converges: its integrand is bounded near zero
and is at most $1/(2t^2)$ for $t\ge1$.  Scaling the integration variable
shows that the first component of the squared distance, if integrated over
all positive scales, would equal $c_*\sqrt m\,|\theta-\eta|$.
The omitted scales contribute at most
\[
 2\sqrt m\,a+\frac{\sqrt m\,|\theta-\eta|^2}{b}
 \le 2\sqrt m\,a+\frac{36m^{3/2}}{b}.
\]
The squared distance in all the additional components is at most
$4\varepsilon SN\log(b/a)$.  Consequently the error is at most
\[
 2m^{-79/2}+36m^{-77/2}+320m^{-191}\log m,
\]
which is at most $m^{-37}$ for sufficiently large $m$.
\end{proof}

\subsection{Extensions with small coordinate profiles}

The auxiliary components have amplitude $\sqrt\varepsilon$, while the
preliminary extension below has amplitude $\varepsilon^{-1/2}$. Their
inner products therefore retain full strength despite the negligible
contribution to squared distances. We use this to extend the derivative
of $p$ from each chart subspace with small coordinate profiles against
all the common features. The extension must agree exactly with $Dp$ along
the chart subspace; its small pairings against $p(\eta)$ will control the
triangle inequalities after rounding. We keep separate polynomial bounds
on ordinary Hilbert norms to control the rounding errors.

\begin{lemma}[Chart derivatives]
\label{ker:extension}
For every $s$ there is a smooth map
$A_s:\R^m\to\mathcal L(\R^N,\Hh)$ such that
\begin{equation}
\label{ker:extension-identity}
 A_s(\theta)J_s=Dp(\theta).
\end{equation}
For $|\theta|,|\eta|\le3\sqrt m$ and $1\le i\le N$,
\begin{equation}
\label{ker:profile}
 \bigl|\langle A_s(\theta)e_i,p(\eta)\rangle\bigr|
 \le C\frac{m}{N}\log m.
\end{equation}
Moreover, uniformly in $\theta\in\R^m$ and $s$,
\begin{equation}
\label{ker:strong-bounds}
 |p(\theta)|,\quad \|Dp(\theta)\|_{\op},\quad
 \|D^2p(\theta)\|_{\op},\quad
 \|A_s(\theta)\|_{\op},\quad
 \|DA_s(\theta)\|_{\op}\le m^{300}.
\end{equation}
Here the norm of a second derivative or an operator-valued derivative is
the corresponding multilinear operator norm.
\end{lemma}

\begin{proof}
We first identify the profile that the preliminary extension should
reproduce. For $|\theta|,|\eta|\le3\sqrt m$, put
$w=(\eta-\theta)/\sigma$. Differentiation in $\theta$ of the Gaussian
term in \eqref{eq:kernel} gives the vector
\[
 \sqrt m\int_a^b w e^{-|w|^2/2}\,\frac{d\sigma}{\sigma}.
\]
The empirical approximation in \Cref{lem:directions} suggests replacing
$w e^{-|w|^2/2}$ by $N^{-1}\sum_i g_{s,i}\sin\langle g_{s,i},w\rangle$.
We realize this profile using the auxiliary Fourier components.
Define $A_s^0(\theta)e_i$ to be supported in summand $(s,i)$ of $\Hh$,
where it is the function
\[
 \frac{m}{N}\varepsilon^{-1/2}
 E'\!\left(\frac{\langle g_{s,i},\theta\rangle}{\sigma}\right).
\]
Its profile is
\begin{equation}
\label{ker:initial-profile}
 \langle A_s^0(\theta)e_i,p(\eta)\rangle
 =\frac{m}{N}\int_a^b\sin\langle g_{s,i},w\rangle\,
          \frac{d\sigma}{\sigma}.
\end{equation}
Since $u_{s,i}=g_{s,i}/\sqrt m$, the profile of
$A_s^0J_s$ is represented by the vector
\[
 \sqrt m\int_a^b
 \left(\frac1N\sum_{i=1}^N g_{s,i}\sin\langle g_{s,i},w\rangle\right)
 \frac{d\sigma}{\sigma}.
\]
Thus $A_s^0J_s$ has the intended empirical profile. Set
$R_s=Dp-A_s^0J_s$; we now bound its pairings against the common features
before correcting the extension on the chart subspace.
For the parameters under consideration,
$|w|\le6\sqrt m/a<m^{42}$.  Thus \Cref{lem:directions} bounds the
difference between these vectors by $2\log(b/a)$.
The gradient of the second term of \eqref{eq:kernel} has norm at most
\[
 \varepsilon\sum_{t=1}^{S}\sum_{i=1}^{N}|g_{t,i}|
        \int_a^b\frac{d\sigma}{\sigma^2}
 \le\frac{2\varepsilon SN\sqrt m}{a}\le1.
\]
It follows that
\begin{equation}
\label{ker:residual-profile}
 \bigl|\langle R_s(\theta)v,p(\eta)\rangle\bigr|
 \le C(\log m)|v|.
\end{equation}
Now define
\begin{equation}
\label{ker:A-definition}
 A_s=A_s^0+R_sJ_s^+=A_s^0(I-J_sJ_s^+)+Dp\,J_s^+.
\end{equation}
The identity $J_s^+J_s=I$ proves \eqref{ker:extension-identity}.
Combining \eqref{ker:initial-profile}, \eqref{ker:residual-profile}, and
the column bound $|J_s^+e_i|\le4m/N$ from \eqref{eq:pseudoinverse}
proves \eqref{ker:profile}.

For completeness, the required Hilbert norm bounds have ample polynomial
slack.  Differentiation under the Hilbert space formulas is valid to every
order, because $a>0$ and Gaussian moments of every order are finite.
For $j=0,1,2$, the squared multilinear derivative norm of the first
component of $p$ is at most
\[
 \sqrt m\,b\,a^{-2j}\E|G|^{2j},
\]
and that of all the additional components together is at most
\[
 \varepsilon SN\log(b/a)
       \left(\frac{2\sqrt m}{a}\right)^{2j}.
\]
Using $\E|G|^2=m$ and $\E|G|^4=m(m+2)$, these estimates give
$\|D^jp\|_{\op}\le m^{110}$ for $0\le j\le2$.
Bounding an operator by $\sqrt N$ times its largest column norm gives
\[
 \|A_s^0\|_{\op}
 \le \sqrt N\,\frac{m}{N}\varepsilon^{-1/2}\sqrt{\log(b/a)},
 \qquad
 \|DA_s^0\|_{\op}
 \le \sqrt N\,\frac{m}{N}\varepsilon^{-1/2}
          \sqrt{\log(b/a)}\frac{2\sqrt m}{a}.
\]
In particular, these two quantities are at most $m^{110}$ and $m^{150}$,
respectively.  Substitute these bounds in
\eqref{ker:A-definition}, together with
$R_s=Dp-A_s^0J_s$, its derivative, and
$\|J_s\|_{\op}\le2\sqrt N$ and
$\|J_s^+\|_{\op}\le2\sqrt{m/N}$.
The projection $I-J_sJ_s^+$ has norm at most one; the last expression
in \eqref{ker:A-definition} and its derivative therefore bound all five
quantities in \eqref{ker:strong-bounds} by $m^{300}$ for sufficiently
large $m$.
\end{proof}

\subsection{Rounded charts and their angles}

Let $Q_s=I-J_sJ_s^+$, the orthogonal projection perpendicular to the
image of $J_s$.  Define
\begin{equation}
\label{ker:P-definition}
 P_s(x)=p(J_s^+x)+A_s(J_s^+x)Q_sx,
 \qquad x\in\R^N.
\end{equation}
On the chart subspace this gives $P_s(J_s\theta)=p(\theta)$ and
$DP_s(J_s\theta)=A_s(\theta)$: use $Q_sJ_s=0$ and
\eqref{ker:extension-identity}. We use only values very close to that
subspace; the derivative away from it is computed explicitly below.
Let $U=(-2,2)^m$, and for $\theta\in U$ put
\begin{equation}
\label{ker:rounding}
 X_s(\theta)_i
 =\tau\left\lfloor\frac{\langle u_{s,i},\theta\rangle}{\tau}
       \right\rfloor.
\end{equation}
In each chart, discard the threshold hyperplanes
$\langle u_{s,i},\theta\rangle=k\tau$, $k\in\mathbb Z$;
their union has Lebesgue measure zero.  Define the finite set
\begin{equation}
\label{ker:vertices}
 V=\{(s,X_s(\theta)):1\le s\le S,\ 
       \theta\in U\text{ off the threshold hyperplanes of chart }s\}.
\end{equation}
Write $v_s(\theta)=(s,X_s(\theta))$ whenever this is defined, and
$P_v=P_s(x)$ for $v=(s,x)\in V$.

For any Hilbert vectors $Z_v$, their squared distances satisfy all
triangle inequalities exactly when
\begin{equation}\label{rep:eq:angle}
\langle Z_v-Z_z,Z_y-Z_z\rangle\ge0
\qquad(v,y,z\in V).
\end{equation}
We next bound how far the vectors $P_v$ can violate this condition.

\begin{lemma}[Two angle estimates]
\label{ker:angles}
For sufficiently large $m$, all $v,z,y\in V$ satisfy
\begin{equation}
\label{ker:global-angle}
 \langle P_v-P_z,P_y-P_z\rangle\ge-r.
\end{equation}
There is an absolute constant $C_0$ such that, with $L=C_0\log m$, for
$v=(s,x)$ and $z=(s,x')$ in the same chart and every $y\in V$,
\begin{equation}
\label{ker:local-angle}
 \bigl|\langle P_v-P_z,P_y-P_z\rangle\bigr|
 \le L\frac{m}{N}\|x-x'\|_1.
\end{equation}
Moreover, for all defined $v_s(\theta),v_t(\eta)$,
\begin{equation}
\label{ker:rounded-distance}
 \left|\,|P_{v_s(\theta)}-P_{v_t(\eta)}|^2
       -c_*\sqrt m\,|\theta-\eta|\,\right|
 \le m^{-37}+m^{610}\tau\le r/10.
\end{equation}
\end{lemma}

\begin{proof}
For $x=X_s(\theta)$, rounding and \eqref{eq:pseudoinverse} give
\begin{equation}
\label{ker:rounding-errors}
 |x-J_s\theta|\le\sqrt N\tau,\qquad
 |J_s^+x-\theta|\le2\sqrt m\tau,\qquad
 |Q_sx|\le\sqrt N\tau.
\end{equation}
In particular, $|J_s^+x|\le3\sqrt m$.  By
\eqref{ker:strong-bounds} and \eqref{ker:P-definition},
\begin{equation}
\label{ker:feature-error}
 |P_{v_s(\theta)}-p(\theta)|
 \le m^{300}(2\sqrt m+\sqrt N)\tau
 \le m^{305}\tau.
\end{equation}
The change in a squared distance caused by the two errors in
\eqref{ker:feature-error} is at most $m^{610}\tau$, since
$|p(\theta)|\le m^{300}$.  Therefore \Cref{lem:kernel} gives
\eqref{ker:rounded-distance}; its last inequality follows directly from
\eqref{ker:parameters}.

Choose parameter representatives in $U$ for $v,z,y$.  Apply
\eqref{ker:rounded-distance} to their three pairs and use
\[
 2\langle P_v-P_z,P_y-P_z\rangle
 =|P_v-P_z|^2+|P_y-P_z|^2-|P_v-P_y|^2.
\]
The corresponding expression for the ordinary Euclidean distances of
the three representatives is nonnegative by the triangle inequality.
The total error is at most $3r/10$, proving
\eqref{ker:global-angle}.

For the local estimate, we integrate coordinate profiles so that the bound
scales with $\|x-x'\|_1$. Consider the segment joining $x$ to $x'$.
Every point $\xi$ on this segment satisfies
$|Q_s\xi|\le\sqrt N\tau$ and $|J_s^+\xi|\le3\sqrt m$, by convexity
and \eqref{ker:rounding-errors}.  Differentiating
\eqref{ker:P-definition} and using
\eqref{ker:extension-identity} gives the exact identity
\begin{equation}
\label{ker:DP-identity}
 DP_s(\xi)h
 =A_s(J_s^+\xi)h
  +DA_s(J_s^+\xi)[J_s^+h]Q_s\xi.
\end{equation}
The operator norm of the second term is at most
$2m^{300}\sqrt m\tau\le m^{305}\tau$.
For a vertex $y=v_t(\eta)$, replace $P_y$ by $p(\eta)$ using
\eqref{ker:feature-error}.  Equations \eqref{ker:profile} and
\eqref{ker:strong-bounds} then imply, throughout this segment,
\[
 \bigl|\langle DP_s(\xi)e_i,P_y\rangle\bigr|
 \le C\frac{m}{N}\log m+3m^{605}\tau
 \le C'\frac{m}{N}\log m.
\]
The same bound holds with $P_z$ in place of $P_y$.
Integrating \eqref{ker:DP-identity} along the segment and summing the
absolute coordinate increments proves \eqref{ker:local-angle} after
choosing a sufficiently large absolute $C_0$.
\end{proof}
